% =====================================================================
\chapter{Valor absoluto y desigualdades}
\label{tema:2}
% =====================================================================

Gran parte del trabajo matemático en aprendizaje automático consiste en \emph{acotar}: garantizar que un error es pequeño, que una pérdida no es negativa o que un método converge. Este tema reúne las desigualdades elementales que se usan para ello. El valor absoluto y la desigualdad triangular siguen el capítulo~1 de Spivak~\cite{spivak}; las inecuaciones, el §2.7 de OpenStax~\cite{openstax}; las demostraciones por inducción, el capítulo~10 de Hammack~\cite{hammack}.

\section{Valor absoluto}

\begin{definicion}[Valor absoluto]\label{def:abs}
Para $x\in\R$ se define $\abs{x}=x$ si $x\ge 0$ y $\abs{x}=-x$ si $x<0$. La \emph{distancia} entre $x$ e $y$ es $\abs{x-y}$.
\end{definicion}

\begin{proposicion}\label{prop:abs}
Para todos $x,y\in\R$ y $r\ge 0$:
\begin{enumerate}
  \item $\abs{x}\ge 0$, con igualdad si y solo si $x=0$; $\abs{-x}=\abs{x}$;
  \item $\abs{x}=\max(x,-x)$ y $-\abs{x}\le x\le\abs{x}$;
  \item $\abs{x}=\sqrt{x^2}$ y $\abs{xy}=\abs{x}\abs{y}$;
  \item $\abs{x}\le r$ si y solo si $-r\le x\le r$; $\abs{x}\ge r$ si y solo si $x\le -r$ o $x\ge r$.
\end{enumerate}
\end{proposicion}

\begin{proof}
(i) y (ii) se comprueban separando los casos $x\ge0$ y $x<0$. (iii) $\abs{x}$ es el único número no negativo cuyo cuadrado es $x^2$, luego coincide con $\sqrt{x^2}$; entonces $\abs{xy}=\sqrt{x^2y^2}=\sqrt{x^2}\sqrt{y^2}=\abs{x}\abs{y}$. (iv) Si $x\ge0$, $\abs{x}\le r$ equivale a $x\le r$, y $-r\le x$ es automático; si $x<0$, $\abs{x}\le r$ equivale a $-x\le r$, es decir, $x\ge -r$, y $x\le r$ es automático. La segunda equivalencia es la negación de la primera con desigualdad estricta.
\end{proof}

\section{La desigualdad triangular}

\begin{teorema}[Desigualdad triangular]\label{teo:triangular}
Para todos $x,y\in\R$, $\abs{x+y}\le\abs{x}+\abs{y}$, con igualdad si y solo si $xy\ge0$.
\end{teorema}

\begin{proof}
Ambos miembros son no negativos, así que basta comparar sus cuadrados (\cref{eje:1-potencia-n}):
\[
(\abs{x}+\abs{y})^2-\abs{x+y}^2=x^2+2\abs{x}\abs{y}+y^2-(x^2+2xy+y^2)=2(\abs{xy}-xy)\ge0,
\]
porque $\abs{xy}\ge xy$ por la \cref{prop:abs}(ii). La diferencia es cero si y solo si $\abs{xy}=xy$, es decir, si y solo si $xy\ge0$.
\end{proof}

\begin{corolario}[Desigualdad triangular inversa]\label{cor:triangular-inversa}
Para todos $x,y\in\R$, $\bigl|\abs{x}-\abs{y}\bigr|\le\abs{x-y}$.
\end{corolario}

\begin{proof}
Por el \cref{teo:triangular}, $\abs{x}=\abs{(x-y)+y}\le\abs{x-y}+\abs{y}$, luego $\abs{x}-\abs{y}\le\abs{x-y}$. Intercambiando los papeles, $\abs{y}-\abs{x}\le\abs{y-x}=\abs{x-y}$. Las dos desigualdades juntas dicen que $\abs{x}-\abs{y}\in[-\abs{x-y},\abs{x-y}]$, que es el enunciado por la \cref{prop:abs}(iv).
\end{proof}

La desigualdad triangular se generaliza a $n$ sumandos, $\abs{\sum_i x_i}\le\sum_i\abs{x_i}$ (\cref{eje:2-triangular-n}), y a vectores, donde $\abs{\cdot}$ se sustituye por una norma (módulo~2).

\section{Inecuaciones con valor absoluto}

Por la \cref{prop:abs}(iv), la condición $\abs{x-a}<r$ describe el intervalo $(a-r,a+r)$ de centro $a$ y radio $r$. Las inecuaciones con varios valores absolutos se resuelven separando casos según el signo de cada expresión.

\begin{ejemplo}\label{ej:abs-casos}
Resolver $\abs{x-1}+\abs{x+2}\ge 5$. Los puntos de corte son $-2$ y $1$.
Si $x\ge1$, la inecuación es $(x-1)+(x+2)\ge5$, es decir, $x\ge2$.
Si $-2\le x<1$, es $(1-x)+(x+2)=3\ge5$, imposible.
Si $x<-2$, es $(1-x)-(x+2)\ge5$, es decir, $x\le-3$.
La solución es $(-\infty,-3]\cup[2,\infty)$: los puntos cuya suma de distancias a $1$ y a $-2$ es al menos~5.
\end{ejemplo}

\section{Desigualdad entre las medias aritmética y geométrica}

\begin{proposicion}[Caso de dos números]\label{prop:amgm2}
Si $a,b\ge0$, entonces $\sqrt{ab}\le\dfrac{a+b}{2}$, con igualdad si y solo si $a=b$.
\end{proposicion}

\begin{proof}
$\dfrac{a+b}{2}-\sqrt{ab}=\dfrac{(\sqrt a-\sqrt b)^2}{2}\ge0$, y la igualdad equivale a $\sqrt a=\sqrt b$.
\end{proof}

\begin{lema}\label{lem:producto1}
Si $x_1,\dots,x_n>0$ y $x_1x_2\cdots x_n=1$, entonces $x_1+\dots+x_n\ge n$, con igualdad si y solo si todos valen~1.
\end{lema}

\begin{proof}
Por inducción sobre $n$. Para $n=1$, $x_1=1$. Supóngase cierto para $n$ y sean $x_1,\dots,x_{n+1}>0$ de producto~1. Si todos valen~1, la suma es $n+1$. Si no, hay alguno mayor que~1 y alguno menor que~1; reordenando, $x_n>1>x_{n+1}$. Entonces $(x_n-1)(1-x_{n+1})>0$, que desarrollado da
\[
x_n+x_{n+1}>1+x_nx_{n+1}.
\]
Los $n$ números $x_1,\dots,x_{n-1},\,x_nx_{n+1}$ tienen producto~1, luego, por la hipótesis de inducción, $x_1+\dots+x_{n-1}+x_nx_{n+1}\ge n$. Sumando,
\[
x_1+\dots+x_{n+1}>x_1+\dots+x_{n-1}+x_nx_{n+1}+1\ge n+1.\qedhere
\]
\end{proof}

\begin{teorema}[Desigualdad de las medias]\label{teo:amgm}
Si $a_1,\dots,a_n>0$, entonces
\[
\sqrt[n]{a_1a_2\cdots a_n}\le\frac{a_1+a_2+\dots+a_n}{n},
\]
con igualdad si y solo si $a_1=a_2=\dots=a_n$.
\end{teorema}

\begin{proof}
Sea $G=\sqrt[n]{a_1\cdots a_n}>0$ y $x_i=a_i/G$. Entonces $x_1\cdots x_n=(a_1\cdots a_n)/G^n=1$ y, por el \cref{lem:producto1}, $\sum_i a_i/G\ge n$, es decir, $\frac1n\sum_ia_i\ge G$. La igualdad se da si y solo si todos los $x_i$ valen~1, es decir, si todos los $a_i$ valen~$G$.
\end{proof}

\begin{ejemplo}\label{ej:amgm-optim}
Para $x>0$, $x+\dfrac 9x\ge 2\sqrt{x\cdot\dfrac9x}=6$, con igualdad si y solo si $x=9/x$, es decir, $x=3$. Así se obtiene un mínimo sin derivar.
\end{ejemplo}

\section{Desigualdades de Bernoulli y de Cauchy--Schwarz}

\begin{teorema}[Desigualdad de Bernoulli]\label{teo:bernoulli}
Si $x\ge-1$ y $n\in\N$, entonces $(1+x)^n\ge1+nx$.
\end{teorema}

\begin{proof}
Por inducción sobre $n$. Para $n=1$ hay igualdad. Si $(1+x)^n\ge1+nx$, como $1+x\ge0$ se puede multiplicar sin cambiar el sentido:
\[
(1+x)^{n+1}\ge(1+nx)(1+x)=1+(n+1)x+nx^2\ge1+(n+1)x.\qedhere
\]
\end{proof}

\begin{teorema}[Desigualdad de Cauchy--Schwarz]\label{teo:cs}
Para reales $a_1,\dots,a_n$ y $b_1,\dots,b_n$,
\[
\Bigl(\sum_{i=1}^na_ib_i\Bigr)^2\le\Bigl(\sum_{i=1}^na_i^2\Bigr)\Bigl(\sum_{i=1}^nb_i^2\Bigr).
\]
\end{teorema}

\begin{proof}
Si todos los $a_i$ son nulos, ambos miembros valen~0. En otro caso, sea $A=\sum a_i^2>0$, $B=\sum a_ib_i$ y $C=\sum b_i^2$. Para todo $t\in\R$,
\[
0\le\sum_{i=1}^n(ta_i+b_i)^2=At^2+2Bt+C.
\]
En particular, para $t=-B/A$ se obtiene $0\le B^2/A-2B^2/A+C=C-B^2/A$, y multiplicando por $A>0$, $B^2\le AC$, que es la desigualdad buscada. (Equivalentemente: el discriminante $4B^2-4AC$ de esa cuadrática no negativa es $\le0$; véase el \cref{eje:3-cuadratica-positiva}.)
\end{proof}

\begin{observacion}[Errores frecuentes]
\begin{enumerate}[label=(\alph*)]
\item $\abs{a+b}=\abs a+\abs b$ es falso en general: $\abs{3+(-5)}=2<8$.
\item La desigualdad de las medias exige números no negativos: con $a=b=-1$, $\sqrt{ab}=1>-1=\frac{a+b}2$.
\item Elevar al cuadrado ambos miembros de una desigualdad solo es válido si ambos son no negativos: $-3<2$ pero $9>4$.
\item En $\abs{x}<r$ con $r\le0$ no hay solución; en $\abs{x}>r$ con $r<0$ lo es todo $\R$.
\end{enumerate}
\end{observacion}

\begin{observacion}[Conexión con el curso]
La desigualdad triangular para normas acota errores acumulados (módulos~7 y~17). Cauchy--Schwarz garantiza que la similitud coseno está en $[-1,1]$ (módulo~2). La desigualdad de las medias dice que la media geométrica de las probabilidades que un modelo asigna a un texto nunca supera su media aritmética; la perplejidad del módulo~1 es precisamente el inverso de esa media geométrica. Su generalización, la desigualdad de Jensen, da la desigualdad de Gibbs del módulo~6.
\end{observacion}

\begin{lecturas}
\item Spivak~\cite{spivak}, cap.~1 (valor absoluto y desigualdad triangular) y sus problemas finales, que incluyen la desigualdad de las medias.
\item OpenStax \emph{College Algebra}~\cite{openstax}, §2.7 «Linear Inequalities and Absolute Value Inequalities».
\item Hammack~\cite{hammack}, cap.~10 «Mathematical Induction», para las demostraciones de desigualdades por inducción.
\end{lecturas}

\section*{Ejercicios del tema 2}
\addcontentsline{toc}{section}{Ejercicios del tema 2}

\begin{ejercicio}{1}{eje:2-calcular}
Calcule $\abs{-3}-\abs{2-7}$ y escriba sin barras $\abs{\pi-4}$ y $\abs{x-1}$ cuando $x<1$.
\end{ejercicio}

\begin{ejercicio}{1}{eje:2-vf}
Razone si son verdaderas o falsas para todos $a,b\in\R$: (a) $\abs{a-b}\le\abs a-\abs b$; (b) $\abs{a+b}=\abs a+\abs b$ si y solo si $ab\ge0$; (c) $\sqrt{x^2}=x$.
\end{ejercicio}

\begin{ejercicio}{1}{eje:2-traducir}
Exprese con un valor absoluto: (a) el intervalo $(2,8)$; (b) el intervalo $[-1,4]$; (c) «$x$ está a distancia menor que $0{,}01$ de~$3$».
\end{ejercicio}

\begin{ejercicio}{2}{eje:2-inec-simple}
Resuelva: (a) $\abs{3x-2}<7$; (b) $\abs{2x-1}=\abs{x+3}$.
\fuente{adaptado de OpenStax \emph{College Algebra}~\cite{openstax}, §2.7.}
\end{ejercicio}

\begin{ejercicio}{2}{eje:2-inec-casos}
Resuelva $\abs{x+1}\ge 2x-4$.
\end{ejercicio}

\begin{ejercicio}{2}{eje:2-inec-suma}
Resuelva $\abs{x-1}+\abs{x+2}\le5$ y compare el resultado con el \cref{ej:abs-casos}.
\fuente{adaptado de OpenStax \emph{College Algebra}~\cite{openstax}, §2.7.}
\end{ejercicio}

\begin{ejercicio}{2}{eje:2-amgm-aplic}
Usando la desigualdad de las medias: (a) halle el mínimo de $x+\frac4x$ para $x>0$; (b) entre todos los rectángulos de perímetro~20, halle el de área máxima; (c) halle el mínimo de $x+y+z$ si $x,y,z>0$ y $xyz=8$.
\end{ejercicio}

\begin{ejercicio}{2}{eje:2-bernoulli-num}
Use la desigualdad de Bernoulli para demostrar, sin calcular las potencias, que $1{,}01^{100}\ge2$ y que $0{,}99^{50}\ge\frac12$. Compare con los valores reales.
\end{ejercicio}

\begin{ejercicio}{2}{eje:2-cs-num}
(a) Compruebe la desigualdad de Cauchy--Schwarz para $a=(1,2,3)$ y $b=(4,-1,2)$. (b) Halle el valor máximo de $2x+3y$ sobre los puntos que cumplen $x^2+y^2=13$ y el punto en que se alcanza.
\end{ejercicio}

\begin{ejercicio}{3}{eje:2-triangular-n}
Demuestre por inducción que $\abs{x_1+\dots+x_n}\le\abs{x_1}+\dots+\abs{x_n}$ para todo $n\in\N$.
\fuente{adaptado de Spivak~\cite{spivak}, cap.~2, problemas finales (inducción).}
\end{ejercicio}

\begin{ejercicio}{3}{eje:2-maxmin}
Demuestre que, para todos $x,y\in\R$,
\[
\max(x,y)=\frac{x+y+\abs{x-y}}{2}\qquad\text{y}\qquad\min(x,y)=\frac{x+y-\abs{x-y}}{2}.
\]
\fuente{adaptado de Spivak~\cite{spivak}, cap.~1, problemas finales.}
\end{ejercicio}

\begin{ejercicio}{3}{eje:2-cadena-medias}
Demuestre que si $0<a<b$, entonces $a<\sqrt{ab}<\dfrac{a+b}{2}<b$.
\fuente{adaptado de Spivak~\cite{spivak}, cap.~1, problemas finales.}
\end{ejercicio}

\begin{ejercicio}{3}{eje:2-media-armonica}
Sean $a,b>0$. Demuestre que
\[
\frac{2ab}{a+b}\le\sqrt{ab}\le\frac{a+b}{2}\le\sqrt{\frac{a^2+b^2}{2}},
\]
(medias armónica, geométrica, aritmética y cuadrática) y que cada igualdad se da si y solo si $a=b$.
\end{ejercicio}

\begin{ejercicio}{3}{eje:2-bernoulli-fuerte}
Demuestre por inducción que, si $x\ge0$ y $n\in\N$, entonces $(1+x)^n\ge1+nx+\dfrac{n(n-1)}{2}x^2$. Deduzca que $2^n>\dfrac{n(n-1)}2$ para todo $n\in\N$.
\fuente{adaptado de Hammack~\cite{hammack}, cap.~10 (desigualdades por inducción).}
\end{ejercicio}

\begin{ejercicio}{3}{eje:2-suma-inversos}
Demuestre que si $x_1,\dots,x_n>0$, entonces $\displaystyle\Bigl(\sum_{i=1}^nx_i\Bigr)\Bigl(\sum_{i=1}^n\frac1{x_i}\Bigr)\ge n^2$, y que la igualdad se da si y solo si todos los $x_i$ son iguales.
\end{ejercicio}
